\section{Cómo se usa el resultado en un soporte cuya talla crece}
Sea $U$ un soporte de $5\le v\le L\log n$ sitios y $c$ su número de componentes.
Introducimos $d_i^*=\max_{f\ni i}|U\cap f|$ para medir el número de vecinos.
Se distingue de $d_i$, número de factores incidentes. Sea $\beta$ el número de
ciclos independientes del grafo bipartito sitios--fibras, incluyendo RC.
Para el potencial entre líneas $V$ con las sustracciones del archivo de grados,
escribimos $[U]V$ para su coeficiente por diferencia alternada en las activaciones.
\paragraph{Premisas de la implicación.}
En los discos $|s|\le\kappa/v$, $|\zeta_i|\le\kappa'/d_i^*$, se necesita un logaritmo
aproximado de norma $O(v)$ para cada cara y el perfil físico $|\zeta_i|\le C/j$.
Se usan los controles de \emph{Grados ponderados}, partes 02--04. En particular
$\sum_{i\in f}1/d_i^*\le1$. Con la norma $\max_i d_i^*|x_i|$ se mantiene acotada
la masa en cada fibra, y la derivada direccional del gradiente de la acción controla
$\sum_l|\partial_l\partial_i R|/d_l^*$ sin multiplicar por el número de vecinos.
En los átomos se usa el balanceo del menor con las medidas de tamaño $m$.
Las diferencias mixtas de coste $l+1$ del peso logarítmico están acotadas por
$(l+1)! C^{l+1}|s|^l$ para $l\ge1$, bajo la representación holomorfa acotada
$s^{-1}\mathscr A+\mathscr B_q$ de la parte 03 citada, en una bola fija de medidas.
Se integra derivada a derivada en las $l+1$ direcciones de supresión, de norma $O(|s|)$.
El criterio de inserción de la parte 04 citada da entonces $|Z|,|1/Z|\le e^{Cv}$
al reducir $\kappa,\kappa'$. La parte 05 de esta entrega controla el error físico
en ese polidisco de \emph{campos}, aun sin prolongación en $s$ del error exacto.
Las premisas analíticas son distintas de las identidades algebraicas de las partes 06--07.
\paragraph{Proyección.}
El exterior común permite aplicar $\prod_{i\in U}(I-\operatorname{ev}_{\zeta_i=0})$
después de la diferencia en las activaciones. Expandir la serie de potencias conserva
sólo monomios con todos los sitios presentes. Evaluar en la actividad física cuesta
$\prod_i(C'd_i^*/j)$ por Cauchy y la suma geométrica en cada sitio. Si un sitio
tiene potencia al menos dos, sumar sobre él cuesta además $\sum_i C'd_i^*/j\le Cv^2/j$.
En la variable $s$, la suma desde orden $w$ cuesta $2(Cv/j)^w$ con $s=1/j$.
En un bosque se toma $w=c+1$ por la parte 06. Para varios ciclos basta $w=c-1$.
Con un solo ciclo el término $w=c-1$ sólo sobrevive con una actividad repetida;
desde $w=c$ basta la ganancia ordinaria. En el caso $c=1$ se usa también la
exactitud del bosque sin reacción y el cálculo del ciclo de coste cero.
Para el error físico, en cada cara la suma sobre átomos cuesta a lo sumo $e^{Cv}\Delta_q$;
la alternancia cuesta a lo sumo $2^v$. Desde cierto $n$, $Cq\le\sqrt{j}$ y
$q/2=2v+6>v+c+1$. Esto absorbe ese error incluso antes de aprovechar la proyección
en campos. En conjunto se recupera la cota de \emph{Grados ponderados}, parte 05,
\[
 |[U]V|\le C^v(Cv)^{c+1}\left(\prod_{i\in U}d_i^*\right)j^{-p},\quad
 p=\begin{cases}v+c+1&\beta=0,\\v+c&\beta=1,\\v+c-1&\beta\ge2.\end{cases}
\]
Es una cota por soporte, condicionada a los márgenes y mayorantes que se acaban de
enumerar. Para la velocidad se necesitan además la suma y la cola sobre soportes,
la clase de llegada y los controles probabilísticos de la identidad detenida.
